% BHU PYQ -- Professional Exam Template
\documentclass[11pt,a4paper]{article}

\usepackage[a4paper,top=2.5cm,bottom=2.5cm,left=2.8cm,right=2.8cm,headheight=14pt]{geometry}
\usepackage{amsmath,amssymb}
\usepackage[T1]{fontenc}
\usepackage[utf8]{inputenc}
\usepackage{lmodern}
\usepackage{microtype}
\usepackage{fancyhdr}
\usepackage{enumitem}
\usepackage{listings}
\usepackage{hyperref}
\usepackage{lastpage}
\usepackage{parskip}

\hypersetup{colorlinks=true,urlcolor=black,linkcolor=black,
  pdftitle={CSB-504: Operating System Concepts -- BHU PYQ}}

\pagestyle{fancy}
\fancyhf{}
\renewcommand{\headrulewidth}{0.4pt}
\renewcommand{\footrulewidth}{0.4pt}
\fancyhead[L]{\small\textit{Banaras Hindu University}}
\fancyhead[R]{\small\textit{CSB-504: Operating System Concepts}}
\fancyfoot[C]{\small Page \thepage\ of \pageref{LastPage}}

\lstset{
  basicstyle=\small\ttfamily,
  frame=single,
  breaklines=true,
  columns=flexible,
  keepspaces=true
}

\newlist{parts}{enumerate}{2}
\setlist[parts,1]{label=(\alph*),leftmargin=*,itemsep=4pt,topsep=3pt}
\setlist[parts,2]{label=(\roman*),leftmargin=*,itemsep=2pt,topsep=2pt}
\newcommand{\pts}[1]{\hfill\textit{[#1]}}

\begin{document}

\begin{center}
  {\large\bfseries BANARAS HINDU UNIVERSITY}\\[2pt]
  {\normalsize B.Sc. (Hons.) Semester V Examination, 2013-14}\\[6pt]
  \rule{0.8\linewidth}{0.5pt}\\[6pt]
  {\large\bfseries Paper: CSB-504 --- Operating System Concepts}\\[4pt]
  {\normalsize Time: Three Hours \hfill Maximum Marks: 70}
\end{center}

\rule{\linewidth}{0.4pt}

\smallskip
\noindent\textit{Instructions: Answer five questions in all, including Question No.1, which is compulsory. Figures on the right-hand side}

\bigskip

Paper: CSB-504 (Operating System Concepts)
Note: Answer five questions in all, including Question No.1, which is compulsory. Figures on the right-hand side

\medskip\noindent\textbf{Question 1.} Answer all the parts
\begin{parts}
  \item Briefly explain the main design goals of an operating System.\pts{3}
  \item What are the differences a process and a thread? Describe some benefits\pts{4}
  using threads.
  \item Explain the difference between internal and external fragmentation. Why\pts{4}
  should they be avoided?
  \item List the four conditions for the occurrence of a deadlock and describe the\pts{5}
  resource-allocation graph.
  \item Briefly discuss the terms: Dynamic Loading, Busy form of waiting, and\pts{6}
  Monitor.
\end{parts}

\medskip\noindent\textbf{Question 2.} (a) Briefly explain the different process states and transition among the states\pts{2}
with transition diagram.
\begin{parts}
  \item Explain the terms Process Contro! Block and Context Switching.\pts{4}
  \item There are five processes A to E to run. Their arrival times are 0,1,3,9 and 12 (6}
  seconds, respectively. Their processing times are 3, 5, 2, 5 and 5 seconds,
  respectively. What is average turn-around times using First-Come-First-
  Served, Shortest-Job-First and Round-Robin (with 1 second Time Quantum)
  Scheduling?
\end{parts}

\medskip\noindent\textbf{Question 3.} (a) Discuss Paging Memory Management Scheme with an illustrative example {8)
and explain the hardware support required for Paging.
\begin{parts}
  \item Consider a paging system with the page table stored in memory.\pts{4}
  \item tfa memory reference takes 200 nanoseconds, how long does a paged
  memory reference take?
  \begin{parts}
    \item if we add TLBs and 75 percent of all page-table references are found in
    TLBs, what is the effective memory reference time? {assume that
    finding a page-table entry in TLBs is 10 times faster than that of the
    memory)
  \end{parts}
\end{parts}

\medskip\noindent\textbf{Question 4.} (a) What is Demand Paging? Explain the steps involved in Page fault handling.\pts{4}
\begin{parts}
  \item Briefly describe Belady’s anomaly. |\pts{2}
  \item A small computer has 4 frames. A process makes the following list of page\pts{6}
  references 1, 2, 3, 4, 1,5, 2, 3, 1, 2. How nany page faults occur using FIFO,
  Second Chance and Least-Recently Used page replacement algorithms?
\end{parts}

\medskip\noindent\textbf{Question 5.} (a) Describe the physical and logical views of a disk storage system.\pts{3}
\begin{parts}
  \item Consider the following sequence of disk track requests: 27, 129, 110, 186,\pts{9}
  147, 41, 10, 64 and 120. Assume that initially the head is at track 30 and is
  moving in the direction of decreasing track number. Compute the number of
  PTO,
  tracks the head traverses using FIFO, SSF and Elevator algorithms.
\end{parts}

\medskip\noindent\textbf{Question 6.} (a) Explain the contiguous and linked file allocation methods with major and the\pts{6}
FAT scheme for linked allocation.
\begin{parts}
  \item Explain the concepts involved in maintaining the file system security {2)
  \item Consider a swapping system in which memory consist of the following hole\pts{4}
  sizes in memory order: 10K, 4K, 20K, 18K, 7K, 12K, and 15K. What hole is
  taken for successive segment requests of
  \item 12K
  \begin{parts}
    \item 10K
    \item 9K
    Using First-fit, best-fit, worst-fit and next-fit strategies.
    7 Write notes on any three of the followings 3x4=12
  \end{parts}
  \item Disk free space management schemes .
  \item Types of Operating System
  \item Semaphore
  \item Deadlock Recovery
  \item 
\end{parts}

\vfill
\rule{\linewidth}{0.4pt}
\begin{center}\small
\href{https://bhu-exam.vercel.app/}{Click here to visit our website}\\[4pt]\href{https://me-aryan.vercel.app/}{Click here to visit our contributor}\\[4pt]\href{https://quashan-me.vercel.app/}{Click here to visit our contributor}
\end{center}

\end{document}
